\chapter{Phobia}

\section*{Definition}
Phobia is an intense, persistent, and disproportionate fear of a specific object, situation, or activity, leading to immediate anxiety or avoidance. Specific phobia is diagnosed when the fear is present for at least 6 months and causes significant distress or functional impairment.

\section*{What Happens in Your Body}
Phobias involve hyperactivation of the amygdala, insula, and anterior cingulate cortex upon exposure to the feared stimulus, with reduced inhibitory input from the ventromedial prefrontal cortex. Impaired fear extinction learning reflects deficits in prefrontal--amygdala circuitry, particularly the inability to safely update threat associations.

\section*{Causes}
\begin{itemize}
  \item \textbf{Psychiatric:} Specific phobia (DSM-5; subtypes: animal, natural environment, blood-injection-injury, situational, other), social anxiety disorder, agoraphobia
  \item \textbf{Genetic:} Heritability of specific phobias ~30--40\%, higher for blood-injection-injury type
  \item \textbf{Learning mechanisms:} Direct conditioning (traumatic encounter), vicarious learning (observing others' fear), informational transmission (warnings, media)
  \item \textbf{Evolutionary:} Preparedness theory -- humans are biologically predisposed to fear stimuli that posed ancestral threats (snakes, spiders, heights, enclosed spaces)
\end{itemize}

\section*{When to See a Doctor}
See your doctor if:
\begin{itemize}
  \item Phobic avoidance significantly restricts daily activities, work, or social interactions
  \item The fear has persisted for more than 6 months and causes marked distress
  \item Individual desires treatment to improve quality of life
\end{itemize}

\section*{Related Symptoms}
\begin{itemize}
  \item Panic attacks (in feared situations), avoidance behaviors, anticipatory anxiety, hypervigilance, social withdrawal, panic disorder, agoraphobia
\end{itemize}
\textbf{Important:} Specific phobias are highly treatable with exposure-based therapies, yet most affected individuals never seek treatment.

\section*{Self-Care / Home Management}
\begin{itemize}
  \item Create a fear hierarchy: list feared situations from least to most frightening
  \item Practice gradual self-exposure, staying in the situation until anxiety decreases
  \item Learn diaphragmatic breathing and applied relaxation for blood-injection-injury phobia
  \item Avoid safety behaviors (e.g., always having an escape route) as they maintain fear
\end{itemize}

\section*{How Your Doctor Will Evaluate You}
\begin{itemize}
  \item Clinical interview using DSM-5 diagnostic criteria
  \item Differentiate specific phobia from social anxiety disorder (fear of negative evaluation) and agoraphobia (fear of inability to escape)
  \item Assess severity and impairment using the Fear Questionnaire or SPIN (Social Phobia Inventory)
  \item Rule out PTSD if the fear stems from a traumatic event
\end{itemize}

\section*{Treatment Options}
\begin{itemize}
  \item First-line: Exposure therapy -- graduated, systematic desensitization, or flooding (therapist-assisted)
  \item Cognitive therapy: restructuring catastrophic beliefs about the feared stimulus
  \item One-session treatment (OST) for specific phobias: massed exposure in a single session
  \item Pharmacotherapy: beta-blockers (blood-injection-injury phobia) or benzodiazepines for occasional use; no medication is first-line for specific phobias
  \item Virtual reality exposure for phobias where in vivo exposure is impractical (flying, heights)
\end{itemize}

\section*{References}
\begin{thebibliography}{99}
\bibitem{phobia:ref1} Bandelow B, Michaelis S, Wedekind D. "Treatment of anxiety disorders." \textit{Dialogues in Clinical Neuroscience}, 2017;19(2):93--107.
\bibitem{phobia:ref2} Choy Y, Fyer AJ, Lipsitz JD. "Treatment of specific phobia in adults." \textit{Clinical Psychology Review}, 2007;27(3):266--286.
\bibitem{phobia:ref3} Wolitzky-Taylor KB, Horowitz JD, Powers MB, et al. "Psychological approaches in the treatment of specific phobias: a meta-analysis." \textit{Clinical Psychology Review}, 2008;28(6):1021--1037.
\bibitem{phobia:ref4} \"Ost LG. "One-session treatment for specific phobias." \textit{Behaviour Research and Therapy}, 1989;27(1):1--7.
\bibitem{phobia:ref5} Curtis GC, Magee WJ, Eaton WW, et al. "Specific fears and phobias: epidemiology and classification." \textit{British Journal of Psychiatry}, 1998;173:212--217.
\bibitem{phobia:ref6} Kaczkurkin AN, Foa EB. "Cognitive-behavioral therapy for anxiety disorders: an update on the empirical evidence." \textit{Dialogues in Clinical Neuroscience}, 2015;17(3):337--346.

\end{thebibliography}
